\hypertarget{independent-rows-without-acousto-optic-frequency-offsets}{%
\chapter{Independent rows without acousto-optic frequency
offsets}\label{row_coherence-independent-rows-without-acousto-optic-frequency-offsets}}

Follow-up to \texttt{docs/stacked\_relay.md} (Sec. 5) and
\texttt{docs/collimator.md} (Sec. 8), which proposed one acousto-optic
frequency shifter per chip. Code: \texttt{scripts/m15\_coherence.py}
(vector focusing in \texttt{picsim/opa/stack.py},
\texttt{Focus2D.field\_vector}); results \texttt{results/m15\_*.json}.
Values are SIM or DERIVED, valid in the regime stated, and should be
reviewed before they are quoted.

\hypertarget{1-result}{%
\section{Result}\label{row_coherence-1-result}}

Per-chip frequency offsets are not needed. The same decoupling is
obtained with one electro-optic phase modulator on the seed laser,
driven with Gaussian noise of about 3 rad rms and 1 to 2 GHz bandwidth,
ahead of the 1:16 split. Each chip\textquotesingle s feed is then given
a different fibre length, in steps of at least three correlation lengths
(about 0.3 m at 2 GHz, at most 5 m in total). Within one chip all paths
agree to better than a millimetre, so the light stays coherent and route
E is unaffected. Between chips the relative phase becomes a random
process with correlation time below 1 ns, so the atoms see the time
average.

A fully passive alternative, orthogonal linear polarisation on alternate
chips, meets the dynamic criteria. It leaves a slow depth wander of
0.4\% rms and needs a polarisation-insensitive objective.

Perturbation of a tweezer by all other rows within +/-16 um (both sides;
rms over drifting relative phases, pairs in quadrature; 4 um rows, full
rows of 41 tweezers):

{\footnotesize
\begin{longtable}[]{@{}>{\raggedright\arraybackslash}p{0.192\dimexpr(\linewidth-10\tabcolsep)\relax}>{\raggedright\arraybackslash}p{0.256\dimexpr(\linewidth-10\tabcolsep)\relax}>{\raggedright\arraybackslash}p{0.081\dimexpr(\linewidth-10\tabcolsep)\relax}>{\raggedright\arraybackslash}p{0.169\dimexpr(\linewidth-10\tabcolsep)\relax}>{\raggedright\arraybackslash}p{0.302\dimexpr(\linewidth-10\tabcolsep)\relax}@{}}
\toprule\noalign{}
option & added hardware & power & static depth / position (rms) & when a
tweezer moves 0.25 um: depth / position (rms) \\
\midrule\noalign{}
\endhead
\bottomrule\noalign{}
\endlastfoot
all chips coherent, no measure & none & 1 & 2.3\% / 9.2 nm & 2.3e-3 /
5.1 nm \\
per-chip AOM offsets & 16 AOMs at 2.6 to 3.4 W & 1 (+ AOM loss) & 1.8e-4
/ 0.09 nm & 2e-5 / 0.05 nm \\
\textbf{common phase noise + fibre delays} & \textbf{1 EOM on the seed,
noise source, \textless= 5 m fibre} & \textbf{1} & \textbf{1.6e-4 / 0.07
nm} & \textbf{2e-5 / 0.04 nm} \\
alternating x / y polarisation & half-wave-plate strips (passive) & 1 &
4.2e-3 / 2.5 nm & 1.1e-3 / 1.6 nm \\
y apodiser a = 0.65 & graded mirror (passive) & 1.84 & 5.2e-3 / 4.1 nm &
1.0e-3 / 1.8 nm \\
polarisation + apodiser a = 0.8 & passive & 1.55 & 2.8e-3 / 1.8 nm &
5.7e-4 / 0.7 nm \\
\end{longtable}
}

(Criteria used for the rows within a chip: others\textquotesingle{}
depth change \textless{} 3e-3 and position change \textless{} 5 nm when
a tweezer moves; proposal budget 10 nm rms position noise.) Leaving all
chips coherent breaks the position budget: 9.2 nm rms of slow wander
plus 5.1 nm when neighbours move.

\hypertarget{2-what-the-frequency-offset-actually-does}{%
\section{What the frequency offset actually
does}\label{row_coherence-2-what-the-frequency-offset-actually-does}}

Light from chip m reaches row m+1 at about 1e-4 of the peak intensity
(module 13). If the two chips are mutually coherent, this field
interferes with the neighbour\textquotesingle s own field, and the depth
changes by about 2 sqrt(epsilon) cos(phi), where phi is the relative
optical phase of the two feeds. Separate fibres let phi drift with
temperature and acoustics. A frequency offset delta f makes phi advance
at 2 pi delta f. When delta f is far above the trap frequencies, the
atom sees only the time average, and the interference term vanishes. Any
process that randomises phi on time scales much shorter than the trap
period does the same. A frequency offset is one such process; a common
random phase seen with different delays is another.

\hypertarget{3-common-phase-modulation-with-fibre-delays}{%
\section{Common phase modulation with fibre
delays}\label{row_coherence-3-common-phase-modulation-with-fibre-delays}}

If the seed carries a phase phi(t) and chip m receives it after a delay
tau\_m, the relative phase of chips m and n is phi(t - tau\_m) - phi(t -
tau\_n). Its mean contribution to the interference is the field
coherence \textbar\textless exp(i {[}phi(t) - phi(t -
tau){]})\textgreater\textbar{} at tau = tau\_m - tau\_n.

{\footnotesize
\begin{longtable}[]{@{}>{\raggedright\arraybackslash}p{0.186\dimexpr(\linewidth-8\tabcolsep)\relax}>{\raggedright\arraybackslash}p{0.291\dimexpr(\linewidth-8\tabcolsep)\relax}>{\raggedright\arraybackslash}p{0.314\dimexpr(\linewidth-8\tabcolsep)\relax}>{\raggedright\arraybackslash}p{0.209\dimexpr(\linewidth-8\tabcolsep)\relax}@{}}
\toprule\noalign{}
modulation & residual coherence, delays 2 to 15 bits & with finite
modulator bandwidth (10-90\% rise 0.3 bit) & at arbitrary (non-integer)
delays \\
\midrule\noalign{}
\endhead
\bottomrule\noalign{}
\endlastfoot
binary 0/pi m-sequence, N = 1023 & 9.8e-4 (= 1/N) & 4.1e-2 (carrier
0.12) & as integer \\
same, depth error +5\% / +10\% & 5.2e-3 / 2.4e-2 & & \\
Zadoff-Chu polyphase, N = 1021 & \textless{} 1e-13 & 4.7e-3 & 6.2e-2 \\
\textbf{Gaussian phase noise, 3 rad rms} & 4.2e-3 (record-limited;
ensemble exp(-9) = 1.2e-4) & inherent & 7.2e-3 \\
Gaussian phase noise, 4 rad rms & 5.6e-3 (record-limited) & inherent &
8.2e-3 \\
\end{longtable}
}

Binary modulation is fragile. Every 0-to-pi and pi-to-0 transition
passes through +pi/2, so a band-limited modulator leaves a carrier (0.12
at 0.3-bit rise), and the delays must be exact multiples of the bit
length. Polyphase codes have perfect autocorrelation but also need exact
delays. Gaussian phase noise of 3 rad rms needs neither: the carrier is
exp(-sigma\^{}2/2), and any delay longer than about three correlation
times decorrelates. The 4e-3 to 8e-3 above is the statistical floor of a
32,768-bit record. What the atom sees is set by the low-frequency
density of the interference term, 0.40/B. For the coherent term of
module 13 (2 sqrt(epsilon) = 2.1e-2), this gives S\_I(2 nu\_trap) =
8.8e-14/Hz at B = 2 GHz. The parametric heating rate Gamma = pi\^{}2
nu\^{}2 S\_I(2 nu) (O\textquotesingle Hara et al.,
arXiv:physics/0003049, used in module 3) is then 6e-3 s\^{}-1 at nu = 84
kHz, an energy e-folding time above 2 minutes. At B = 0.5 GHz it is
2.4e-2 s\^{}-1.

\textbf{Spectral width.} With 3 rad of noise at correlation rate B, the
optical spectrum is Gaussian with rms width 0.67 B (1.35 GHz, 5 pm, at B
= 2 GHz). Its effect on the tweezers comes through the diffractive
objective\textquotesingle s axial chromatic shift (16 um/nm, DERIVED for
f = 17.3 mm with f proportional to 1/lambda) and the lateral scaling x
proportional to lambda:

{\small
\begin{longtable}[]{@{}>{\raggedright\arraybackslash}p{0.092\dimexpr(\linewidth-10\tabcolsep)\relax}>{\raggedright\arraybackslash}p{0.218\dimexpr(\linewidth-10\tabcolsep)\relax}>{\raggedright\arraybackslash}p{0.172\dimexpr(\linewidth-10\tabcolsep)\relax}>{\raggedright\arraybackslash}p{0.218\dimexpr(\linewidth-10\tabcolsep)\relax}>{\raggedright\arraybackslash}p{0.299\dimexpr(\linewidth-10\tabcolsep)\relax}@{}}
\toprule\noalign{}
B & optical rms width & axial rms smear & peak & lateral rms smear at 60
um \\
\midrule\noalign{}
\endhead
\bottomrule\noalign{}
\endlastfoot
0.5 GHz & 0.34 GHz & 0.02 um & 0.99997 & 0.07 nm \\
2 GHz & 1.35 GHz & 0.08 um & 0.9998 & 0.29 nm \\
10 GHz & 6.7 GHz & 0.41 um & 0.995 & 1.4 nm \\
\end{longtable}
}

The spectrum is symmetric about the carrier, so the mean light shift at
the double-magic wavelength is unchanged to first order. The
second-order term needs the curvature of the differential polarisability
{[}VALUE?{]}.

\textbf{Coherence inside a chip.} Path differences within a chip and its
relay are at most about 1 mm (about 5 ps). With a 0.5 ns correlation
time (B = 2 GHz) and sigma = 3 rad, the coherence loss is sigma\^{}2
(tau/tau\_c)\^{}2 / 2 \textasciitilde{} 5e-4. The least-squares
synthesis of route E therefore sees an essentially coherent source.

\textbf{Hardware.} A fibre-coupled LiNbO3 phase modulator on the
low-power seed and a Gaussian noise source amplified to about 1 V\_pi
rms (3 rad) over 1 to 2 GHz. Extra fibre is 0, 0.3, 0.6 ... 4.5 m on the
16 chip feeds, before or after the amplifiers. A seed broadened to about
1 GHz also raises the stimulated-Brillouin threshold of the high-power
amplifier and delivery fibres, which carry 3 W each {[}REF?, RECALLED:
SBS suppression by phase modulation is standard in high-power
single-frequency fibre amplifiers{]}. If the chosen amplifier already
broadens its seed for this reason, the only addition is the fibre
delays. Compared with 16 AOMs, this removes 16 RF drivers and the AOM
insertion loss (about 15\% each {[}VALUE?{]}).

\hypertarget{4-passive-options}{%
\section{Passive options}\label{row_coherence-4-passive-options}}

\textbf{Orthogonal polarisation on alternate chips.} A half-wave-plate
strip in front of every other chip, in the collimated region after the
collimators, rotates its light to y. We computed vector focal fields
(Richards-Wolf, aplanatic NA 0.5 objective) and summed the intensity
over Ex, Ey and Ez (scalar light shift):

{\small
\begin{longtable}[]{@{}>{\raggedright\arraybackslash}p{0.222\dimexpr(\linewidth-8\tabcolsep)\relax}>{\raggedright\arraybackslash}p{0.286\dimexpr(\linewidth-8\tabcolsep)\relax}>{\raggedright\arraybackslash}p{0.302\dimexpr(\linewidth-8\tabcolsep)\relax}>{\raggedright\arraybackslash}p{0.190\dimexpr(\linewidth-8\tabcolsep)\relax}@{}}
\toprule\noalign{}
separation & polarisations & static: depth / position (worst phase) &
moving: depth / position \\
\midrule\noalign{}
\endhead
\bottomrule\noalign{}
\endlastfoot
4 um & orthogonal & 8.4e-4 / 1.2 nm & 1.6e-4 / 0.3 nm \\
8 um & parallel & 3.3e-3 / 2.1 nm & 1.1e-3 / 1.4 nm \\
12 um & orthogonal & 5.0e-5 / 0.3 nm & 4.4e-5 / 0.1 nm \\
16 um & parallel & 2.4e-3 / 0.4 nm & 3.5e-4 / 0.6 nm \\
\end{longtable}
}

At NA 0.5 the orthogonal neighbours\textquotesingle{} interference falls
by a factor of 26 (the residual comes from the depolarisation of
high-angle rays, mainly via Ez). The rows two apart (8 um) share a
polarisation and then dominate. The moving-tweezer criteria are met, and
the position budget too (2.5 nm rms). The static depth wanders by 0.4\%
rms with the fibre phases, at Hz to kHz rates, far below twice the trap
frequencies.

Three costs follow. (i) The meta-objective must be
polarisation-insensitive (propagation-phase meta-atoms with fourfold
symmetry rather than geometric phase). (ii) Where the orthogonal fields
overlap, the polarisation becomes elliptical in proportion to their
relative phase, which produces a vector light shift. Our estimate, from
the far-detuned alkali formula for Cs at 1064 nm, is a vector-to-scalar
ratio of about 0.09 g\_F m\_F per unit ellipticity (RECALLED, verify).
With an ellipticity of up to about 2e-2 at 4 um, this is about 2e-3 g\_F
m\_F of the trap depth: up to tens of kHz for stretched states, zero to
first order for clock states (m\_F = 0). (iii) Only two orthogonal
states exist, so rows two apart stay coherent.

\textbf{Passive y apodiser.} A Gaussian graded-transmission (or
graded-reflectivity) filter at the plane of the y telescope conjugate to
the pupil apodises all chips at once. It lowers the 1D crosstalk at 4 um
from 8.4e-5 to 3.5e-6 (a = 0.65 pupil radii, beta = 0.7 um). However, it
costs a factor 1.84 in power and widens the y spot from 1.35 to 1.73 um.
In 2D, the coherent sum of a full row, with the circular pupil, still
leaves 4.7e-3 of static depth change at 4 um. It is not competitive.

\hypertarget{5-recommendation}{%
\section{Recommendation}\label{row_coherence-5-recommendation}}

Use common phase modulation with fibre delays in place of per-chip AOMs:
one modulator, no power penalty, and the same row independence as
frequency offsets. Alternating polarisation costs nothing to add, and it
gives a second, independent factor of about 26 on the nearest rows if
the modulation fails or is switched off.

Next step (i) is simulated in \texttt{docs/pm\_am.md} (filters narrower
than about 100 pm and reflection pairs with R1 R2 \textgreater{} 1.7e-3
must be avoided; a representative chain gives -134.5 dB/Hz). Next steps:
(i) measure the delayed self-coherence of the modulated seed after the
amplifier (PM-to-AM conversion by any spectral filtering adds intensity
noise at GHz; its density at 2 nu\_trap should be checked); (ii) choose
the noise drive (or a deterministic multi-tone drive with the same
autocorrelation); (iii) obtain the second-order differential light shift
of the 1064.4350 nm double-magic condition over +/-5 pm.

\hypertarget{verification-notes}{%
\section{Verification notes}\label{row_coherence-verification-notes}}

{\footnotesize
\begin{longtable}[]{@{}>{\raggedright\arraybackslash}p{0.262\dimexpr(\linewidth-8\tabcolsep)\relax}>{\raggedright\arraybackslash}p{0.131\dimexpr(\linewidth-8\tabcolsep)\relax}>{\raggedright\arraybackslash}p{0.179\dimexpr(\linewidth-8\tabcolsep)\relax}>{\raggedright\arraybackslash}p{0.429\dimexpr(\linewidth-8\tabcolsep)\relax}@{}}
\toprule\noalign{}
value & location & type & pointer \\
\midrule\noalign{}
\endhead
\bottomrule\noalign{}
\endlastfoot
option table & Sec. 1 & SIM/DERIVED & results/m15\_summary.json from
m15\_coherent and m15\_prbs (review) \\
2D coherent / incoherent, vector fields & Sec. 1, 4 & SIM &
results/m15\_coherent.json (review) \\
residual coherence, spectra, chromatic smear & Sec. 3 & SIM &
results/m15\_prbs.json (review) \\
apodiser scan & Sec. 4 & SIM & results/m15\_apod.json (review) \\
heating rates & Sec. 3 & DERIVED & O\textquotesingle Hara et al.,
arXiv:physics/0003049 (used and checked in module 3), density from
m15\_prbs \\
16 um/nm axial chromatic shift & Sec. 3 & DERIVED & f/lambda for a
diffractive objective, f = 17.3 mm \\
vector-to-scalar ratio 0.09 g\_F m\_F & Sec. 4 & RECALLED & far-detuned
formula (Grimm et al., arXiv:physics/9902072), not read this session;
counter-rotating terms neglected \\
3 rad rms, B 1-2 GHz, 0.3 m delay steps & Sec. 1, 3 & CHOICE & design
choices \\
\end{longtable}
}

Open placeholders: {[}REF?{]} SBS suppression by phase modulation in
high-power 1064 nm fibre amplifiers; {[}VALUE?{]} AOM insertion loss at
3 W; {[}VALUE?{]} second-order differential light shift of the
double-magic condition over a 5 pm spectrum; vector polarisability of Cs
at 1064 nm from a primary source.
